\section{Holonomie affine et symétries de la mémoire
}
\label{nuclear:6}
\subsection{Composition effective de la mémoire
}

Pour une émission admissible \(\epsilon\), la construction de \cref{iv:cont:concatenacion} définit


\begin{equation}
\mathsf U_\epsilon(m)=C_\epsilon m+\kappa_\epsilon,
\qquad
\kappa_{\epsilon_2\epsilon_1}
=\kappa_{\epsilon_2}+C_{\epsilon_2}\kappa_{\epsilon_1}.
\label{nuc06:eq:1}
\end{equation}

Le domaine est le module de mémoire de la source ; le codomaine, celui de l’émission suivante. Le produit conserve l’ordre. Sur un chemin \(\gamma\) de longueur \(r\), on obtient

\begin{equation}
C_\gamma=C_r\cdots C_1,\qquad
\kappa_\gamma=\sum_{j=1}^r C_r\cdots C_{j+1}\kappa_j,
\qquad
m_\gamma=C_\gamma m_0+\kappa_\gamma.
\label{nuc06:eq:2}
\end{equation}

Le produit vide de la somme est l’identité. La démonstration est la récurrence de \eqref{nuc06:eq:1} : ajouter l’émission suivante multiplie le registre précédent par son transport linéaire et ajoute son incrément. La frontière positionnelle peut s’annuler tandis que la somme transportée de \eqref{nuc06:eq:2} demeure.


\begin{proposition}[Classe de mémoire d’un retour]

Soit \(M\) un module abélien et soit un chemin fermé dans la base dont le transport de mémoire soit \(H(m)=Cm+\kappa\), avec \(C\) un automorphisme de \(M\). Sa classe de retour est

\begin{equation}
\mathfrak m(H)=[\kappa]\in M/(I-C)M.
\label{nuc06:eq:3}
\end{equation}

Un changement de coordonnées de mémoire \(m'=Rm+b\), avec \(R\) inversible, transforme le retour en

\begin{equation}
C'=RCR^{-1},\qquad
\kappa'=R\kappa+(I-C')b.
\label{nuc06:eq:4}
\end{equation}

L’application induite par \(R\) envoie \([\kappa]\) sur \([\kappa']\). Le transport \(H\) possède un point fixe exactement lorsque sa classe \eqref{nuc06:eq:3} est nulle.
\end{proposition}

\begin{proof}
Substituer \(m=R^{-1}(m'-b)\) dans \(H\) donne \eqref{nuc06:eq:4}. Comme \(R(I-C)=(I-C')R\), \(R\) induit un isomorphisme entre les quotients. Le terme \((I-C')b\) représente zéro dans le quotient d’arrivée. Enfin, l’équation de point fixe \(Cm+\kappa=m\) équivaut à \((I-C)m=\kappa\). Cela prouve la dernière affirmation.

\end{proof}

La formule \eqref{nuc06:eq:3} permet de décider si un changement d’origine élimine l’incrément ou ne fait que le redistribuer entre les coordonnées. En particulier, lorsque \(C=I\), la classe est le vecteur \(\kappa\) lui-même. Une translation non nulle reste non nulle sous tout changement affine inversible.

\begin{proposition}[Symétrie du retour et changement de représentation]

Une transformation affine \(A(m)=Rm+b\) commute avec \(H\) exactement lorsque

\begin{equation}
RC=CR,\qquad (I-C)b=(I-R)\kappa.
\label{nuc06:eq:5}
\end{equation}
\end{proposition}

\begin{proof}
Les parties linéaires de \(AH\) et \(HA\) sont \(RC\) et \(CR\) ; leurs parties constantes sont \(R\kappa+b\) et \(Cb+\kappa\). Les égaler donne \eqref{nuc06:eq:5}. Si l’on transporte seulement \(H\) au moyen de \eqref{nuc06:eq:4}, la commutation n’est pas requise : on obtient la même loi dans une autre représentation. Pour une symétrie du retour fixé, en revanche, les deux égalités de \eqref{nuc06:eq:5} sont requises.
\end{proof}

On distingue ainsi avec précision la covariance de la loi et le stabilisateur d’une mémoire concrète. Cette différence est pertinente pour étudier les variations de symétrie de l’état sans modifier les lois de composition.


\subsection{Application au registre effectif du tour observable}

La section~\ref{vii:sec:retorno-memoria-gravedad} définit deux compteurs ordonnés : \(g\) compte les inversions d’orientation et \(s\) les changements effectifs de feuillet. Leur calcul canonique par bloc de 27 pas est

\[
c_{27}=(1,18).
\]

La phase observable complète revient après quatre de ces blocs. Dans le groupoïde des routes, le module de registre est \(M=\mathbb Z^2\), et


\begin{equation}
c_{108}=(4,72)=4v,\qquad v=(1,18),\qquad
H(g,s)=(g+4,s+72).
\label{nuc06:eq:6}
\end{equation}

L’identité \eqref{nuc06:eq:6} est un résultat du corpus utilisé ici ; le vérificateur vérifie ses conséquences, sans recompter les émissions d’origine. Le symbole \(c_{108}\) désigne cet incrément de deux compteurs. Il est distinct du sceau dodécaphasique \(K\) et du cocycle central de la représentation de Heisenberg.

\begin{theorem}[Invariants complets du registre sous des tours complets]

Les orbites de \(H\) et de son inverse sur \(\mathbb Z^2\) sont classées exactement par


\begin{equation}
\boxed{\quad \mathcal I(g,s)=s-18g,\qquad \nu(g,s)=[g]_4.\quad}
\label{nuc06:eq:7}
\end{equation}

En conséquence,

\begin{equation}
\mathbb Z^2/\mathbb Z(4,72)\cong\mathbb Z\oplus\mathbb Z/4\mathbb Z.
\label{nuc06:eq:8}
\end{equation}
\end{theorem}

\begin{proof}
Ajouter \((4,72)\) conserve les deux expressions de \eqref{nuc06:eq:7}. Si deux registres ont les mêmes invariants, \(g'-g=4N\) pour un entier \(N\) et \(s'-s=18(g'-g)=72N\) ; ils diffèrent donc exactement de \(N(4,72)\). L’application de \eqref{nuc06:eq:7} est surjective : étant donnés \(j\) entier et \(r\) modulo quatre, choisir un représentant \(g\) de \(r\) et poser \(s=j+18g\). Cela démontre \eqref{nuc06:eq:8}.
\end{proof}

La description couvre toutes les orbites de ce lecteur de mémoire. Les orbites de l’état enrichi complet conservent en outre leurs autres coordonnées. Pour un prolongement dirigé vers l’avant, un registre ultérieur s’obtient lorsque l’entier N précédent est positif ou nul ; l’équivalence bilatérale de \eqref{nuc06:eq:8} appartient au groupoïde.


La grandeur \(\mathcal I\) est l’unique charge linéaire entière primitive, au signe près, invariante sous ces translations. En effet, pour \(L(g,s)=ag+bs\), l’invariance exige \(4a+72b=0\), de sorte que \((a,b)=b(-18,1)\). Un tour accumule de la mémoire longitudinale et conserve une relation transversale déterminée par l’incrément.

Le facteur quatre possède ici une origine précise : c’est \(\gcd(4,72)\) et, dans le calcul précédent, le nombre de blocs de 27 qui complète le retour observable. La classe \(\nu\) conserve l’information perdue lorsqu’on divise l’incrément par quatre et qu’on ne garde que la direction primitive \(v\).

\begin{corollary}[Registre conjoint de bloc et de phase
]

Écrivons la phase de ces quatre blocs sous la forme \(a\in\mathbb Z/4\mathbb Z\). La mise à jour


\[
(a,g,s)\longmapsto(a+1,g+1,s+18)
\]

conserve \((s-18g,[g-a]_4)\). Les deux données classent ses orbites bilatérales : la différence \(N=g'-g\) satisfait à la fois \(a'-a=[N]_4\) et \(s'-s=18N\). Cela réunit l’avancée de la mémoire et le retour de la phase sans remplacer l’un par l’autre.

\end{corollary}

\subsection{Générateur parabolique déterminé par la mémoire}

La paire ordonnée de compteurs dote le module \(M\) de la forme alternée unimodulaire

\[
\Omega_M((g,s),(g',s'))=gs'-sg'.
\]

Son domaine est le module de registre. Cette forme n’est identifiée ni à une aire physique dimensionnelle ni à la forme du tore APP, qui a un autre domaine. Le choix de l’ordre \((g,s)\) fixe son orientation.

\begin{theorem}[Stabilisateur unimodulaire de l’incrément]

Définissons, directement à partir du vecteur primitif \(v\) de \eqref{nuc06:eq:6},

\begin{equation}
N_v m=v\,\Omega_M(v,m)=v\mathcal I(m),\qquad
N_v=\begin{pmatrix}-18&1\\-324&18\end{pmatrix}.
\label{nuc06:eq:9}
\end{equation}

Alors \(N_v^2=0\), \(N_v\ne0\) et

\begin{equation}
\boxed{\quad
\operatorname{Stab}_{SL_2(\mathbb Z)}(c_{108})
=\{I+nN_v:n\in\mathbb Z\}.
\quad}
\label{nuc06:eq:10}
\end{equation}

Chaque transformation de \eqref{nuc06:eq:10} conserve \(\Omega_M\), l’incrément \(c_{108}\) et la charge \(\mathcal I\). Pour \(n\) non nul, c’est une matrice unipotente différente de l’identité, de trace deux, appartenant au type parabolique.

\end{theorem}

\begin{proof}
\(\Omega_M(v,v)=0\) implique \(N_v^2m=v\Omega_M(v,v)\Omega_M(v,m)=0\). Son coefficient supérieur droit est un, de sorte qu’il est non nul. La matrice

\[
B=\begin{pmatrix}1&0\\18&1\end{pmatrix}
\]

a pour déterminant un, envoie le premier vecteur de base sur \(v\) et satisfait

\begin{equation}
N_v=BNB^{-1},\qquad N=\begin{pmatrix}0&1\\0&0\end{pmatrix}.
\label{nuc06:eq:11}
\end{equation}

Un automorphisme entier fixe \(4v\) exactement lorsqu’il fixe \(v\), car \(M\) est sans torsion. Une matrice de déterminant un qui fixe le premier vecteur de base a nécessairement la forme \(\begin{pmatrix}1&n\\0&1\end{pmatrix}\). La conjuguer par \(B\) prouve \eqref{nuc06:eq:10}. Le reste se déduit de cette expression ou directement de \eqref{nuc06:eq:9}.

\end{proof}

La formule \eqref{nuc06:eq:9} rend la construction indépendante du second vecteur choisi pour compléter \(v\) en une base orientée unimodulaire : tout autre complément diffère d’un multiple entier de \(v\) et produit le même \(N_v\). Sous une transformation orientée \(R\in SL_2(\mathbb Z)\), on a \(N_{Rv}=RN_vR^{-1}\).

L’action peut s’écrire avec une clarté particulière :


\begin{equation}
(I+nN_v)m=m+n\mathcal I(m)v.
\label{nuc06:eq:12}
\end{equation}

La charge transversale détermine de combien cette symétrie déplace le long de la direction de mémoire. La somme des tours produit le déplacement \(4Nv\) ; les symétries stabilisatrices préservent cette loi d’avancement.

Sur le quotient \eqref{nuc06:eq:8}, la transformation \(I+nN_v\) agit comme \((\mathcal I,[g]_4)\mapsto(\mathcal I,[g+n\mathcal I]_4)\). Par conséquent, elle conserve la collection d’orbites et peut permuter ses classes résiduelles ; \(\nu\) est invariant du retour \(H\), et non de chaque transformation stabilisatrice. Cette différence se déduit directement de \eqref{nuc06:eq:12}.

La conjugaison \eqref{nuc06:eq:11} fournit une application algébrique explicite du modèle parabolique du TRIT, défini dans \cref{euler-trit-generadores}, au module de registre étendu aux scalaires réels. On obtient en outre \(\exp(tN_v)=I+tN_v\) et \(B\exp(tN)=\exp(tN_v)B\). Cet entrelacement concerne des opérateurs : le mouvement particulier sélectionné par le TPK utilise aussi ses émissions et son calendrier. Le groupe arithmétique complet de \eqref{nuc06:eq:10} fixe ce lecteur ; le groupe des symétries admissibles de l’état total s’obtient en prenant son intersection avec les actions qui préservent les autres données et relations.

Dans le système translaté \eqref{nuc06:eq:6}, les changements affines \(m\mapsto Rm+b\) avec \(R\) dans \eqref{nuc06:eq:10} et \(b\) arbitraire commutent avec \(H\). Ceux qui préservent en outre le registre initial fixé doivent satisfaire \(Rm_0+b=m_0\). Distinguer les deux conditions sépare la symétrie de la loi de la symétrie d’une trajectoire désignée.

\begin{theorem}[Réalisation variationnelle de la charge et du stabilisateur
]
\label{nuc06:variacional}

L’extension scalaire du registre, \(M_{\mathbb R}=M\otimes_{\mathbb Z}\mathbb R\), conserve les coordonnées engendrées \((g,s)\) et la forme \(\Omega_M=dg\wedge ds\). C’est une réalisation ultérieure du module entier. En fixant la convention \(\iota_{X_H}\Omega_M=dH\), on a


\begin{equation}
X_{\mathcal I}=v,\qquad
X_{\frac12\mathcal I^2}=\mathcal I v=N_vm.
\label{nuc06:eq:12a}
\end{equation}

Ainsi, la translation et la collection parabolique construites admettent des générateurs hamiltoniens explicites. Le flot de \(\mathcal I\) est \(m\mapsto m+uv\) ; en \(u\) égal à quatre, il reproduit \eqref{nuc06:eq:6}. Le flot de \(\mathcal I^2/2\) est \(m\mapsto(I+uN_v)m\). Le paramètre \(u\) correspond à cette réalisation continue ; le pas discret de retour et les émissions intermédiaires maintiennent leur calendrier TPK propre.
\end{theorem}

\begin{proof}
Pour un vecteur \(X=(X_g,X_s)\), \(\iota_X(dg\wedge ds)=X_g\,ds-X_s\,dg\). Comme \(d\mathcal I=ds-18\,dg\), le premier champ est \((1,18)=v\). La règle de dérivation en chaîne donne le second. \(\mathcal I\) reste constant le long des deux champs, parce que \(d\mathcal I(v)=0\). Par conséquent, le second flot intègre \(dm/du=\mathcal I(m_0)v\), équivalent à \(\exp(uN_v)=I+uN_v\). La non-dégénérescence de \(\Omega_M\) implique que l’hamiltonien qui engendre \(N_v\) est \(\mathcal I^2/2\) plus une constante additive : sa différentielle est fixée par \(\iota_{N_vm}\Omega_M\).

\end{proof}

L’action mathématique

\begin{equation}
\mathscr A[m]=\int_{u_0}^{u_1}
\left[s\frac{dg}{du}-\frac12(s-18g)^2\right]du
\label{nuc06:eq:12b}
\end{equation}

produit par variation les équations \(dg/du=\mathcal I\), \(ds/du=18\mathcal I\). La translation \(g\mapsto g+\varepsilon\), \(s\mapsto s+18\varepsilon\) conserve le Hamiltonien et modifie le lagrangien de la dérivée totale \(d(18\varepsilon g)/du\), pour \(\varepsilon\) constant. Sa charge de Noether est

\begin{equation}
\underbrace{s\,\delta g}_{s}
-\underbrace{18g}_{\text{terme de bord}}
=s-18g=\mathcal I.
\label{nuc06:eq:12c}
\end{equation}

La variation s’interprète par unité de \(\varepsilon\). Les équations vérifient directement \(d\mathcal I/du=18\mathcal I-18\mathcal I=0\). La translation discrète \eqref{nuc06:eq:6} elle-même est également réalisée au moyen du Hamiltonien \(H_{\rm ret}=4\mathcal I\), dont le flux de paramètre un est \(H\) ; les deux Hamiltoniens commutent au sens de Poisson. Le flux stabilisateur \(I+uN_v\) et le retour \(m+4v\) sont des actions distinctes qui commutent : sur la droite \(\mathcal I=0\), le premier fixe chaque point, tandis que le second continue à accumuler de la mémoire.

Plus généralement, tout Hamiltonien \(C^1\) sur ce plan qui est invariant sous les translations \(m\mapsto m+uv\) a la forme \(F(\mathcal I)\). Pour le prouver, les coordonnées globales \((g,\mathcal I)\) écrivent cette symétrie comme une translation de \(g\) ; l’invariance élimine la dépendance en \(g\) de \(H\). Son champ est \(F'(\mathcal I)v\). Cette classification explicite aussi bien la conservation commune que le choix supplémentaire d’une dynamique concrète.

L’action \eqref{nuc06:eq:12b} est adimensionnelle et appartient au plan de registre réalisé. Sa démonstration fournit une relation variationnelle exacte entre mémoire, charge et stabilisateur. La conversion en action physique et en temps dimensionnel se réalise au moyen des lecteurs d’action et d’horloge correspondants ; ces unités ne sont pas attribuées au compteur u.


\subsection{Courbure d’APP et conservation de la phase centrale
}

La section~\ref{iv:extension-curvatura-mixta} travaille sur \(X_9=(\mathbb Z/9\mathbb Z)^2\). Les feuillets sont \(S(x,y)=x+y\) et \(P(x,y)=xy\). Les connexions mixtes utilisent des différences de ces potentiels, et non les potentiels comme liens :

\begin{equation}
F(S,P)=\Delta_x\Delta_yP-\Delta_y\Delta_xS=1,
\qquad F(P,S)=-1\pmod9.
\label{nuc06:eq:13}
\end{equation}

La somme orientée des liens sur le contour d’un rectangle est \(ab\) ou \(-ab\) selon l’ordre. L’aire alternée est \(\sigma(u,v)=ad-bc\) pour \(u=(a,b)\), \(v=(c,d)\). La section~\ref{iv:weyl-area-extension} représente ces transports au moyen de \(D(u)=X^aZ^b\), \(ZX=\omega XZ\), avec \(\omega\) la racine neuvième dans la représentation ultérieure. Alors

\begin{equation}
D(u)D(v)=\omega^{bc}D(u+v),\qquad
D(u)D(v)D(u)^{-1}D(v)^{-1}=\omega^{-\sigma(u,v)}I.
\label{nuc06:eq:14}
\end{equation}

La position résultante du commutateur revient. La phase centrale retient l’aire orientée avec le signe fixé dans \eqref{nuc06:eq:14}. C’est une rétention modulo neuf : une aire multiple de neuf a une phase centrale triviale, bien qu’un relèvement avec quotient et registre puisse distinguer ses parcours. La phase ne remplace pas ce relèvement.


\begin{proposition}[Covariance exacte du transport central]

Comme deux est inversible modulo neuf, définissons

\begin{equation}
W(a,b)=\omega^{ab/2}D(a,b),\qquad
W(u)W(v)=\omega^{-\sigma(u,v)/2}W(u+v).
\label{nuc06:eq:15}
\end{equation}

Pour \(R\in SL_2(\mathbb Z/9\mathbb Z)\), l’application

\begin{equation}
\omega^zW(u)\longmapsto\omega^zW(Ru)
\label{nuc06:eq:16}
\end{equation}

est un automorphisme de l’extension centrale. Il conserve ses produits et ses phases de commutation.
\end{proposition}

\begin{proof}
La phase lors de la multiplication des deux \(W\) est
\(\frac{ab}{2}+\frac{cd}{2}+bc-\frac{(a+c)(b+d)}2=\frac{bc-ad}{2}\). Cela prouve \eqref{nuc06:eq:15}. L’identité \(\sigma(Ru,Rv)=\det(R)\sigma(u,v)=\sigma(u,v)\) montre que \eqref{nuc06:eq:16} conserve la loi ; \(R\) inversible fournit l’automorphisme inverse. Si \(\det R=-1\), l’automorphisme correspondant change aussi \(z\mapsto-z\) ; le signe de l’aire et le centre s’inversent conjointement.
\end{proof}

Le registre intégral et sa phase centrale ont des types différents. La covariance permet de conserver la loi d’aire tandis que la présentation change ; la classe affine de \eqref{nuc06:eq:3} permet de conserver le contenu de mémoire tandis que son origine change. Les deux propriétés agissent sur des lecteurs typés de routes, et leurs applications doivent rester séparées lors de leur composition avec la lecture exceptionnelle de l’état complet.

\begin{proposition}[Persistance sous repères locaux et subdivisions
]

Soit une boucle \(\gamma\) basée en \(x\) et soient ses liens inversibles \(U_\epsilon\). Le changement local \(U'_\epsilon=G_tU_\epsilon G_s^{-1}\) transforme son holonomie en \(H'_\gamma=G_xH_\gamma G_x^{-1}\) : les repères des sommets intérieurs s’annulent successivement. Si \(H_\gamma=\omega^kI\), la phase centrale reste exactement identique. Pour des liens raffinés dont le produit reproduit le lien précédent, la subdivision conserve également \(H_\gamma\), par associativité. Dans les plaquettes APP, l’annulation des arêtes intérieures matérialise cette propriété au moyen de Stokes fini.

\end{proposition}

Un rectangle élémentaire a un Wilson additif égal à un et une phase centrale non triviale. Ce transport mixte distingue donc une holonomie réelle de l’identité que produiraient des liens de la forme \(G_tG_s^{-1}\). La conservation sous subdivision concerne des lacets et des liens compatibles : ajouter une autre aire change le transport.

Le relèvement entier de l’holonomie distingue l’aire orientée de son résidu nonadique : pour un rectangle de côtés quatre et sept, \(28=1+9\cdot3\), tandis que \(-28=8+9(-4)\) correspond à son orientation inverse. La phase résiduelle coïncide pour les aires un et vingt-huit, tandis que leurs quotients distinguent les accumulations. Cet exemple exhibe la différence entre phase finie et registre intégral ; son quotient d’aire et le compteur temporel \(q_{\rm mem}\) conservent des domaines distincts.

\subsection{Composition entre les trois régimes et continuité du bilan
}

Les réalisations de \cref{euler-trit-generadores,euler-trit-evaluaciones} construisent des générateurs \(J_k^2=-\tau_kI\) sur leurs fibres propres, de dimension réelle deux. On prend des liens effectivement déclarés \(B_k:V_k\to V_{k+1}\) qui préservent les formes alternées, \(B_k^*\Omega_{k+1}B_k=\Omega_k\). Le signe \(*\) dans cette section signifie transposition/tiré en arrière bilinéaire. Les opérateurs \(E_k=\exp(t_kJ_k)\) préservent \(\Omega_k\) : en dimension deux, la trace nulle du générateur implique \(J_k^*\Omega_k+\Omega_kJ_k=0\). L’identité symplectique est l’hypothèse qui serait conservée en étendant le lemme à des dimensions supérieures.

Pour la compression nonadique et son complément enregistré,


\[
T_k=B_k(8I+E_k)/9,\qquad
D_k=B_k\sqrt8(E_k-I)/9,
\]

le développement algébrique donne


\begin{equation}
\Omega_k=T_k^*\Omega_{k+1}T_k+D_k^*\Omega_{k+1}D_k.
\label{nuc06:eq:17}
\end{equation}

En effet, les termes croisés s’annulent et le numérateur est \(72\Omega_k+9E_k^*\Omega_kE_k=81\Omega_k\). Pour \(P_0=I\), \(P_{k+1}=T_kP_k\), on obtient, par substitutions successives et dans l’ordre de la route,


\begin{equation}
\boxed{\quad
\Omega_0=P_N^*\Omega_NP_N+
\sum_{k=0}^{N-1}(D_kP_k)^*\Omega_{k+1}(D_kP_k).
\quad}
\label{nuc06:eq:18}
\end{equation}

Le terminal et chaque complément restent présents. L’identité est bilinéaire orientée ; les résultats de positivité s’appliquent dans les profils positifs démontrés dans les sections précédentes. La forme conservée n’impose pas qu’un lien entre régimes distincts entrelace leurs générateurs : si \(BJ_\tau=J_{\tau'}B\), élever au carré avec \(B\) inversible imposerait \(\tau=\tau'\). La conservation d’une structure commune permet de maintenir les trois régimes distincts.

Le transport intégral \(U_N=(B_{N-1}E_{N-1})\cdots(B_0E_0)\) envoie \(V_0\) sur \(V_N\). Une monodromie matricielle requiert une fermeture typée \(K:V_N\to V_0\) ; alors seulement \(KU_N\) est un endomorphisme dont la trace peut être calculée. Des changements de base compatibles conservent cette trace par conjugaison. Retour de phase, fermeture de fibre et retour du vecteur restent des conditions différentes.


Si les liens ont exclusivement la forme \(B_{yx}=G_yG_x^{-1}\), leur produit le long d’un chemin est \(G_{\rm final}G_{\rm inicial}^{-1}\). En refermant avec le même repère, il donne l’identité. L’holonomie de \eqref{nuc06:eq:14} et l’incrément de \eqref{nuc06:eq:6} montrent des mécanismes concrets que cette simple substitution de repères ne peut représenter à elle seule.

\subsection{Compatibilité avec le raffinement et portée arbitraire
}

L’action de mémoire est compatible avec un homomorphisme de modules de restriction \(p:M'\to M\) lorsque


\begin{equation}
pC'=Cp,\qquad p\kappa'=\kappa.
\label{nuc06:eq:19}
\end{equation}

Alors \(pH'=Hp\) et \(p\) induit une application de quotients

\begin{equation}
M'/(I-C')M'\longrightarrow M/(I-C)M,
\qquad [\kappa']\longmapsto[\kappa].
\label{nuc06:eq:20}
\end{equation}

\begin{proof}
La première égalité de \eqref{nuc06:eq:19} envoie \((I-C')M'\) sur \((I-C)M\) ; la seconde fixe l’image de l’incrément. La compatibilité de \(H\) résulte de la substitution. Concaténations et raffinements composent les mêmes carrés grâce à la loi \eqref{nuc06:eq:1}.

\end{proof}

Une classe de retour non nulle qui apparaît à un niveau conservé reste non nulle dans tout relèvement compatible : une classe nulle se projetterait sur zéro. Cette assertion s’applique au système de mémoires et de restrictions qui satisfait \eqref{nuc06:eq:19} ; le constructeur conjoint et les applications de son système projectif sont conservés dans \cref{iv:continuo-conjunto}.

Dans le registre de \eqref{nuc06:eq:6}, pour chaque entier \(N\), on a \(H^N(m)=m+N(4,72)\). La preuve est algébrique pour \(N\) arbitraire, et non une inférence à partir d’un nombre fini d’itérations. Dans le prolongement vers l’avant, chaque préfixe conserve son registre fini ; l’histoire complète est décrite au moyen de la collection compatible de ces registres, sans remplacer ses compteurs par une quantité entière « infinie ».

Si une dynamique complète possède le lecteur \(q_{\rm mem}\) tel que \(q_{\rm mem}(\Gamma_9x)=q_{\rm mem}(x)+1\), elle ne peut pas non plus posséder d’état périodique de période \(N\) positive dans ce domaine : appliquer le lecteur produirait \(q_{\rm mem}(x)=q_{\rm mem}(x)+N\). Ce corollaire utilise le degré entier complet ; ses quotients résiduels peuvent avoir des périodes finies, comme il convient au calendrier observable du corpus.


\subsection{Conséquence fondationnelle et portée de la composition}

Le lien spécifique obtenu est

\begin{equation}
\text{incrément de retour }4v
\ \longrightarrow\ \mathcal I(m)=\Omega_M(v,m)
\ \longrightarrow\ N_vm=v\mathcal I(m)
\ \longrightarrow\ I+nN_v.
\label{nuc06:eq:21}
\end{equation}

Le registre accumulé détermine une charge conservée et une collection parabolique qui stabilise sa direction. L’évolution et la conservation appartiennent à la même action : \(g\) et \(s\) changent, leur relation \(s-18g\) demeure, et l’accroissement conserve un stabilisateur explicite. Cela précise une partie du lien entre mémoire et symétrie qui a motivé la recherche.

L’égalité \eqref{nuc06:eq:18} relie cette recherche aux bilans précédents : une évolution peut redistribuer la forme entre lecture et compléments en conservant le total. Les identités d’APP \eqref{nuc06:eq:13}–\eqref{nuc06:eq:16} montrent, dans un autre lecteur concret de même origine, comment l’ordre du transport produit une donnée centrale même lorsque la position revient.

La charge \eqref{nuc06:eq:7} est un invariant algébrique du transport et, au moyen de l’action explicite \eqref{nuc06:eq:12b}, une charge de Noether de la réalisation du registre. Le théorème~\ref{nuc06:variacional} et son développement variationnel construisent cette relation, sa symétrie et son terme de bord. Les développements antérieurs de l’action et du courant physiques conservent leurs domaines et leurs unités ; les transporter à cette réalisation exige de maintenir leurs applications dimensionnelles.


La composition avec les lecteurs dimensionnels doit transporter la paire \((v,\Omega_M)\) et la charge \(\mathcal I\), en préservant leurs unités et le courant variationnel. Toute relation ultérieure avec \(\pi,\phi,e,\alpha\) et \(K\) doit provenir de cette composition ; les formules présentes fixent l’information à conserver.
